\chapter{Monitoring and controlling reasoning}
\label{sec:p19-rationality-8-6}

Solving a problem, judging how well the solution is proceeding, and deciding
whether to continue are different activities. A person can have a correct
answer with low confidence, an incorrect answer with high confidence, or
an accurate judgment of difficulty while using an ineffective strategy.
\emph{Meta-reasoning} studies the monitoring and control processes associated
with reasoning, problem solving and decision making.

This chapter reconstructs Thompson, Elqayam and Ackerman's
\emph{Reasoning, Rationality, and Metacognition}, Chapter~8.6 of the 2021
\emph{Handbook of Rationality} \cite{ThompsonElqayamAckerman2021MetaReasoning}.
Its account connects metacognitive judgments to epistemic, normative,
practical and bounded rationality. A separate primary evidence record below
uses Stupple and colleagues' 2017 reasoning-time study
\cite{StuppleEtAl2017Slower}. The finite calculations are handbook
constructions. They do not turn this psychological framework into a
universally optimal policy for how long to think.

\section{Judgment, performance and control}
\label{sec:p19-r86-judgments}

\begin{hbdefinition}{def:p19-r86-monitoring-control}{The monitored process and the control action}
The \emph{object-level activity} produces or evaluates an answer, representation
or choice. A \emph{monitoring judgment} concerns that activity: whether a
problem is solvable, whether one can solve it, whether an initial answer feels
right, whether progress is satisfactory, or whether the final answer is correct.
A \emph{control action} changes the activity or its allocation of resources:
continue, reconsider, change strategy, seek information or help, withhold a
response, stop, or commit to an answer.

The framework permits monitoring to be automatic and experienced as a feeling;
it does not require deliberate reflection on every cognitive step. Calling a
process metacognitive does not establish that it is accurate. Its object-level
and monitoring distinction is a psychological distinction; it is not an
identification of cognitive processes with formal object-language and
metalanguage syntax.
\end{hbdefinition}

The source illustrates the distinction with uncertainty about a familiar
person's name and withholding it; initial pessimism about assembling furniture,
followed by an attempt and later help seeking; and dissatisfaction with a
manuscript followed by a break. These are explanatory examples, not experimental
demonstrations that those actions were optimal. A sense that an answer is
unsatisfactory may concern persuasion, accuracy, completeness or a task goal;
the relevant criterion must be specified
\cite[pp.~517--518]{ThompsonElqayamAckerman2021MetaReasoning}.

\begin{hbdefinition}{def:p19-r86-judgment-types}{Judgments at different stages}
An \emph{initial judgment of solvability} concerns a problem before extended
work. Intrinsic solvability and one's own chance of finding a solution are
distinct possible targets. A \emph{feeling of rightness}, abbreviated \(\mathrm{FOR}\),
accompanies an initial answer. \emph{Intermediate confidence} concerns progress
or a current answer during work. \emph{Final confidence} concerns the answer
eventually given. A \emph{feeling of knowing} concerns whether currently
unavailable knowledge could be retrieved or recognized.

A rating scale for any one of these judgments must be recorded. An ordinal
FOR scale is not automatically a probability scale. A numerical confidence
report can be assessed against \ref{def:p19-calibration-law} only when its
predicted event and probability interpretation are specified.
\end{hbdefinition}

Low initial FOR can accompany reconsideration, while low confidence after a
long attempt can accompany stopping or seeking help. These associations are
compatible because the judgments occur at different stages and concern
different control choices. The two-response procedure discussed by the chapter
asks for an initial rapid answer and its FOR, then a considered answer and
final confidence. It records changes and persistence; a longer second phase
does not itself establish more correct responding. The same answer can be
repeated after either verification or unsuccessful reconsideration
\cite[pp.~518, 522]{ThompsonElqayamAckerman2021MetaReasoning}.

\section{Cues to confidence and the argument about rationality}
\label{sec:p19-r86-cues}

The chapter's cue account separates three sources of information: beliefs
about one's own competence, expectations about a task, and features of the
particular item or its processing. Familiarity, the amount of associated
information retrieved, the ease of generating an answer, and consistency
among retrieved considerations can affect a judgment without establishing
the answer's correctness. The proposed connection is a mechanism for
monitoring and control, not direct introspective access to the reliability
of an inference \cite[pp.~518--520]{ThompsonElqayamAckerman2021MetaReasoning}.

\begin{hbdefinition}{def:p19-r86-cue-standing}{A cue, its validity and its use}
A \emph{cue} is a feature available to the reasoner that can affect a
monitoring judgment. \emph{Cue validity} concerns its relation to the actual
target under a specified task/population law. \emph{Cue use} concerns its
relation to the reported judgment or control decision. A cue can be used
strongly while being weakly related to correctness. Conversely, a valid cue
may receive little weight. These are empirical relationships requiring
separate measurements.
\end{hbdefinition}

The source's examples make these distinctions consequential. Someone can
expect to recover a city's name because related facts are accessible without
currently retrieving that name. Familiarity with a toilet or toaster can be
mistaken for an ability to explain its mechanism. In the overclaiming studies
reviewed by the authors, respondents claimed familiarity with invented
financial concepts, including ``fixed-rate deduction''. These are examples
selected by the chapter to illustrate
\emph{meta-ignorance}, not a rule that familiar knowledge is always unreliable.
The chapter also discusses overconfidence among poor CRT performers and
cross-task confidence tendencies. Its review does not supply a common
participant sample, a universal prevalence or a population confidence interval
for these different phenomena
\cite[pp.~518--519]{ThompsonElqayamAckerman2021MetaReasoning}.

Its confidence-to-action illustrations include deciding whether to submit an
answer for scoring and whether to act in a fictitious medical decision task.
The chapter interprets inaccurate confident actions as commission errors and
withheld correct responses as omission errors. Those descriptions concern
their task criteria; deciding which action is preferable also requires the
consequences of error and withholding. No real clinical recommendation follows
from those illustrative tasks
\cite[p.~519]{ThompsonElqayamAckerman2021MetaReasoning}.

Other examples concern coherence and conflict. The reviewed remote-associate
paradigm distinguishes a word triad that has a common associate, such as the
association of struck, beam and light with moon, from a triad without the
specified common associate. The chapter cites above-chance judgments after
a 1.5-second presentation, while also noting qualifications from subsequent work.
That result concerns a defined association task; it does not prove that a
feeling of coherence identifies a valid argument. Likewise, a conflict between
a tempting response and a normatively supported response can lower confidence
without producing the correct response. Detecting conflict and resolving it
are separate outcomes
\cite[pp.~519--520]{ThompsonElqayamAckerman2021MetaReasoning}.

The self-consistency account described in the source concerns choices between
two alternatives. It proposes that agreement among considerations favoring an
alternative helps determine confidence and fluent choice. Agreement among
retrieved reasons can be produced by redundant information or a biased search;
it is not logical consistency of a theory or coherence with the probability
axioms. The chapter's note~3 explicitly distinguishes this psychological
coherence from Bayesian coherence. Its note~4 concerns predictable lexical
associations, rather than truth preservation by a calculus.

\begin{hbconvention}{def:p19-r86-rationality-standards}{Keep the standard of evaluation explicit}
\emph{Normative correctness} evaluates a response under an identified standard,
such as an arithmetic equation or a logical/probabilistic requirement.
\emph{Epistemic rationality}, in the truth-conducive conception used by the
chapter, concerns true belief and reliable knowledge assessment.
\emph{Practical rationality} concerns pursuit of stated goals with their
consequences and costs. \emph{Bounded rationality} additionally attends to
feasible computation, limited time, attention and memory.
These standards can disagree and are not identified by observed frequency.
In the Evans--Over terminology used by the chapter, practical rationality is
``rationality 1'' and normative correctness is ``rationality 2''.
The source's note~2 records disagreement about truth-conduciveness and its
relation to practical success; its working conception does not settle that
philosophical dispute.
\end{hbconvention}

For the bat-and-ball example, if ball price is \(x\) and bat price is
\(x+1\) dollars, then \(2x+1=1.10\) gives \(x=.05\).
Confidence in .10 does not change those equations. But the expected utility
of further checking depends on time, error consequences and available
methods, not solely on whether an initial answer was wrong. Satisficing means
accepting an adequate option under a stopping criterion; it need not mean
accepting an invalid inference. The chapter's contribution is to examine the
psychological cues and changing criteria behind stopping
\cite[pp.~519--521]{ThompsonElqayamAckerman2021MetaReasoning}.

\section{Effort, stopping and persistence}
\label{sec:p19-r86-persistence}

The source separates four claims often grouped into a cognitive-miser
explanation: people stop too soon; they possess the resources needed to finish;
that stopping produces a normative error; and more investment should improve
normative responding. Those claims require different evidence. A participant
can spend little time because the answer is known, or considerable time while
lacking a relevant method. A fast incorrect answer and a slow incorrect answer
do not identify the same process. The source contrasts premature termination
with \emph{labor in vain}: persistence with little prospect of improving the
answer. It does not follow that all deliberation is wasteful or that every
difficult item should be abandoned
\cite[pp.~520--522]{ThompsonElqayamAckerman2021MetaReasoning}.

\begin{hbdefinition}{def:p19-r86-stopping-model}{A diminishing criterion with a time limit}
A schematic operationalization uses an ordered finite set of observation times
\(0=t_0<t_1<\cdots<t_m=T\), an intermediate judgment trajectory
\(q_j\in[0,1]\), and an acceptance criterion \(h_j\in[0,1]\).
The first threshold stop occurs when \(q_j\geq h_j\); if this has
not happened earlier, stop at the time limit \(T\). A diminishing-criterion
model lets the required judgment level decrease as time is invested. A
fixed-criterion model keeps it constant.

The functions here are a handbook representation of the chapter's account,
not fitted equations or universal parameter values from Ackerman's experiments.
Increasing confidence within an item, lowering an acceptance criterion, and
comparing final confidence across items are different operations.
\end{hbdefinition}

\begin{hbtheorem}{thm:p19-r86-stopping}{proved=annex:p19.r86.stopping}{Threshold stopping and a cross-item contrast}
The capped first-threshold rule is well defined on every finite trajectory.
With the same trajectory and time limit, a criterion no larger at any time
cannot produce a later stopping index. It does not determine whether the
selected answer is correct.

For a continuous illustrative trajectory on \([0,10]\), let
\(h(t)=.9-.005t\). The trajectory \(q_E(t)=.6+.04t\)
first meets it at \(t=20/3\), with confidence \(13/15\).
The trajectory \(q_H(t)=.3+.005t\) does not meet it before the
time cap and ends at confidence .35. Both trajectories increase with time,
but the item receiving more time has lower final confidence.
\end{hbtheorem}

The chapter illustrates its diminishing-criterion discussion with an
Experiment~3 comparison from Ackerman (2014): approximately ten seconds and
confidence above 90 percent for easy items, versus about ninety seconds and
confidence near 20 percent for difficult items. This is the 2021 chapter's
reported illustration; it provides neither the original experiment's full
participant record nor interval estimates for those two numbers. It cannot
by itself identify the causal benefit of an additional eighty seconds, since
item difficulty differs. The original study is required for a stronger
task-level empirical assessment
\cite[p.~521]{ThompsonElqayamAckerman2021MetaReasoning}.

The source also discusses continuing information search beyond the first
discriminating cue, sometimes paying for further information. Such behavior
can be relevant to confidence, alternative goals, or the adequacy of the
claimed heuristic model. Calling information redundant requires the actual
decision rule, cue reliability and costs; it cannot be inferred simply from
the number of searches. Its proposed inverted-U relation between resources
and normative responding is a hypothesis suggested by the reviewed findings,
not a universal dose-response law for thinking time
\cite[pp.~521--522]{ThompsonElqayamAckerman2021MetaReasoning}.

\section{A primary record for response time and correctness}
\label{sec:p19-r86-stupple}

\begin{hbdefinition}{def:p19-r86-stupple-record}{The 2017 three-item CRT study}
Stupple and colleagues studied 391 undergraduates: 304 at Lancaster University
and 87 psychology undergraduates at the University of Derby. Participants
already familiar with the task were excluded and replaced; the published
method does not state a total count of those replacements. The final cohort
is a university sample, not a probability sample of all reasoners.
Three answers per person do not create three independently recruited
participants; the participant remains the relevant unit for the reported
overall-score relationship.

The standard bat-and-ball, widget-production and lily-pad-doubling questions
were presented singly by computer in randomized order. Response time ran from
item presentation to submission of the entered answer. Total correct score
ranged from zero to three. Individual answers were classified as the designated
correct answer, the designated intuitive error, or another error. The paper's
names ``analytic'' and ``intuitive'' for response categories do not directly
measure an unobserved reasoning process.

The analysis used \(\log_{10}\) transformations for response-time regressions
and nonparametric comparisons of response categories, with Bonferroni-adjusted
follow-up thresholds. Confidence and FOR were not measured in this experiment.
Its time measure includes reading, calculation and response entry; it is not
a direct quantity of analytic effort. Collection dates, trial-by-trial
strategies and a representative population frame are not supplied by this
method record \cite[pp.~6--7]{StuppleEtAl2017Slower}.
\end{hbdefinition}

The paper's reported association between its transformed aggregate time measure
and total correct score was
\(r=.181,p<.001,N=391\). This was a small observed association, not a
randomized effect of requiring more time. The multiple regression predicting
\emph{total} correct score from the three logged item times gave standardized
coefficients .344 for bat and ball, .073 for widgets, and -.155 for lily pads;
their reported \(p\)-values were below .001, .182 and .004, respectively.
Those coefficients concern an overall score, not each item's separate causal
response to an added second. The paper prints no confidence intervals for
these coefficients \cite[pp.~7--9, table~2]{StuppleEtAl2017Slower}.

The item comparisons preserve a different denominator and question. Median
times in seconds for correct, designated intuitive-error, and other-error
answers were, respectively, with interquartile ranges in parentheses:
\[
\begin{array}{l|rrr}
 \text{Item}&\text{Correct}&\text{Designated error}&\text{Other error}\\\hline
 \text{Bat and ball}&34.97\ (36.64)&19.38\ (14.96)&27.31\ (42.74)\\
 \text{Widgets}&28.81\ (32.38)&24.29\ (20.98)&34.30\ (44.20)\\
 \text{Lily pads}&28.19\ (26.16)&24.95\ (22.92)&50.63\ (53.41)
\end{array}
\]
The corresponding category counts were \((117,259,15)\),
\((119,209,63)\), and \((163,179,49)\), each summing to 391.
The printed interquartile ranges are dispersion summaries, not standard
errors. Correct and designated-error times differed for bat and ball
at the adjusted threshold; the corresponding comparisons for widgets and
lily pads had \(p=.09\) and \(p=.25\). For lily pads, other errors
took longer than both remaining categories, with \(p<.001\) for both
comparisons \cite[pp.~10--12, tables~4--5]{StuppleEtAl2017Slower}.

These observations constrain a pure explanation of poor CRT scores as
unwillingness to spend time. They do not uniquely distinguish missing skills,
an ineffective calculation, rationalization, uncertainty, or individual speed
differences. The authors discuss those alternatives and the limited consistency
of a three-item test. A nonsignificant category comparison is not equivalence,
and an observed slow error does not establish that more thinking caused the
error. The study supplies no confidence-to-action effect because no confidence
judgments or action-withholding manipulation were collected
\cite[pp.~12--15]{StuppleEtAl2017Slower}.

\section{A worked monitoring and control contrast}
\label{sec:p19-r86-worked}

\begin{hbdefinition}{def:p19-r86-trace}{A two-response reasoning record}
For \(n>0\) declared problems, record initial/final answers, their specified
correctness values \(y_i^{(0)},y_i^{(1)}\in\{0,1\}\), confidence
reports \(c_i^{(0)},c_i^{(1)}\in[0,1]\), and extra deliberation time
\(t_i\geq0\). Here confidence predicts the correctness of that phase's
answer; it is not a rating of how hard the item was. Let \(N_{01}\)
count changes from incorrect to correct and \(N_{10}\) changes from
correct to incorrect. Answers can change within the same correctness category,
so these two counts do not exhaust all answer changes.
\end{hbdefinition}

\begin{hbexample}{def:p19-r86-trace-example}{Six fictional problems with declared answer keys}
Use the following answer keys, with classical propositional entailment in
the two logical-inference questions. Problem~1 is the ball-price equation
\(2x+1=1.10\), whose answer is .05. Problem~2 asks whether
\(A\to B,B\) entail \(A\): they do not. Problem~3 asks whether
probabilities .6 and .7 for two events permit union probability .5: they do
not. Problem~4 asks whether \(A\to B,\neg A\) entail \(\neg B\):
they do not. Problem~5 asks whether \(P\lor\neg P\) is true under
every classical Boolean valuation: it is. Problem~6 asks whether
\(m(m-1)\) is even for every integer \(m\): it is.

The following stipulated record is not from either source:
\[
\begin{array}{c|rrrrrr}
 i&1&2&3&4&5&6\\\hline
 y_i^{(0)}&0&1&0&1&1&0\\
 y_i^{(1)}&0&1&1&0&1&0\\
 c_i^{(0)}&.90&.55&.50&.60&.80&.50\\
 c_i^{(1)}&.95&.85&.70&.80&.90&.60\\
 t_i\text{ (seconds)}&20&40&30&35&5&60
\end{array}
\]
For the incorrect entries, let the submitted answers be .10, ``yes'', or
``false'' as appropriate; for correct entries use the keys above. One answer
is corrected and one correct answer is replaced by an error. Correctness
remains three out of six although every reported confidence increases.
\end{hbexample}

\begin{hbtheorem}{thm:p19-r86-trace-calculations}{proved=annex:p19.r86.trace-calculations}{Accuracy, confidence, ranking and time are distinct}
For every record, the accuracy change is \((N_{01}-N_{10})/n\).
In the six-problem record both accuracies are \(1/2\), mean initial
confidence is \(77/120\), and final confidence is \(4/5\).
The extra times sum to 190 seconds. Binary squared confidence loss changes
from \(137/480\) to \(27/80\), worsening despite increased confidence.
Gamma, defined as in \ref{def:p19-jol-gamma} but for correctness of these
reasoning answers, changes from \(1/3\) to \(1/9\).

If final answers are reported only when final confidence is at least .8,
four answers are reported and two are correct. Coverage is \(2/3\),
reported-answer accuracy is \(1/2\), and a correct answer is withheld.
These are exact descriptive calculations on a fictional record, not population
estimates or causal effects of added time.
\end{hbtheorem}

\begin{hbtheorem}{thm:p19-r86-calibrated-control}{proved=annex:p19.r86.calibrated-control}{What calibrated confidence permits in selective reporting}
Under the finite law of \ref{def:p19-calibration-law}, let \(Y=1\)
mean the reported candidate answer is correct, and let \(C\) be calibrated
for that event. For \(h\in[0,1]\), if \(\mathbb P(C\geq h)>0\),
\[
       \mathbb P(Y=1\mid C\geq h)=E(C\mid C\geq h)\geq h.
\]
This does not follow from gamma equal to one, and it requires the same law
for calibration and reporting. It gives no optimal choice of \(h\) without
the consequences of omitted and incorrect answers.
\end{hbtheorem}

The calibration and gamma results are mathematical interfaces that apply when
their variables and law have been specified. The noun-pair experiments in
\ref{def:p19-jol-record} predict later recall under their actual cue and
delay procedure. They do not establish that a delayed reasoning confidence
rating will improve inference, or that a particular confidence threshold will
be calibrated in a new task. Strictly increasing transformations can preserve
gamma while changing numerical calibration, as
\ref{thm:p19-jol-gamma} proves. Empirical transfer needs a new evidence record
under \ref{def:p19-evidence-record}.

\section{The value of another reasoning step}
\label{sec:p19-r86-further-deliberation}

The chapter closes its persistence argument with four considerations from
Douven's account: the adequacy of current options, the anticipated chance of
finding a better option, the cost of search, and the available time. It connects
the first two to FOR and solvability judgments, and uses a diminishing
criterion to emphasize that costs and time limits can interact. This is an
account of further deliberation; knowing confidence in a current answer alone
does not specify the expected improvement from another step
\cite[pp.~522--523]{ThompsonElqayamAckerman2021MetaReasoning}.

\begin{hbdefinition}{def:p19-r86-step-value}{An explicit one-step decision model}
Suppose a specified belief model assigns probability \(p_0\) to a current
answer being correct and \(p_1\) to the answer from a particular additional
procedure being correct. Let a correct answer earn value \(R>0\), an
incorrect answer incur loss \(L\geq0\), and the additional procedure
cost \(\kappa\geq0\). Both options concern the same task and consequences.
The cost can include elapsed time only after a value for that time has been
specified. Feasibility and any hard deadline are separate constraints.
\end{hbdefinition}

\begin{hbtheorem}{thm:p19-r86-step-value}{proved=annex:p19.r86.step-value}{The conditional value of one additional procedure}
Under the stated model, the extra procedure has larger expected net value
exactly when \((R+L)(p_1-p_0)>\kappa\); equality gives indifference.
With \(R=1,L=0,p_0=.6\), an improvement to \(p_1=.9\) at cost .1
is worthwhile under those values, while improvement to .61 at the same cost
is not. A current confidence value does not determine the decision without
the prospective improvement and cost. This is a conditional decision result,
not an empirical assertion that people know those probabilities or costs.
\end{hbtheorem}

\begin{hbmethod}{def:p19-r86-evidence-method}{Evaluate monitoring and control separately}
Specify the task's answer criterion, the judgment's predicted event and scale,
the stage of the judgment, and the available control actions. Record confidence,
correctness and time separately. Assess cue validity and cue use, calibration
and ranking, corrected and harmful changes, and the consequences of withholding
or continuing. Retain participant/item dependence, familiarity, exclusions,
ability/skill measures and the uncertainty of reported comparisons. Use an
intervention design to evaluate the effect of a proposed control policy;
association between duration and performance is not that effect.

State whether the proposed goal is a correct answer, calibrated confidence,
a useful decision, or less time under a specified loss. Preserve contrary
results and task differences. A successful mathematical stopping rule does
not show that an observed psychological criterion implements it. A useful
heuristic cue in one environment need not retain that validity elsewhere.
\end{hbmethod}

\begin{proofboundary}{def:p19-r86-boundary}{Scope of the account}
The 2021 chapter supplies the source-specific framework and argument.
Its literature illustrations have different tasks and populations, and several
reviewed studies are not supplied with full sample/uncertainty records in that
chapter. The bounded primary CRT record supplies an observed time/correctness
comparison, not a general metacognitive mechanism or a randomized effect of
confidence. The worked calculations establish distinctions under their stated
definitions. Neither the review nor those calculations supports a universal
instruction to deliberate longer, stop sooner, trust fluent answers, or transfer
a memory-monitoring intervention to every reasoning task.
\end{proofboundary}
